\section{Exponential normalization for Gamma and Barnes functions}
\label{sec:exponential-normalization}

Power--logarithmic scales describe algebraic growth and its corrections.
A function such as $\Gamma(x)^2$ grows faster than every power of $x$.
Its logarithm, however, has a power--logarithmic expansion. The appropriate
construction retains the nonvanishing part of that logarithm exactly and
expands only the vanishing part. This yields a relative asymptotic expansion
whose absolute error includes the extracted growth factor.

\subsection{From absolute logarithmic error to relative error}

Let $u\to0^+$ and write $\M(u)=1+|\log u|$. Suppose a real function has fixed
sign $\sigma\in\{1,-1\}$ sufficiently near the endpoint and
\begin{equation}
 \log|f(u)|=A(u)+B(u)+R(u),\qquad
 R(u)=\OO\bigl(u^\rho \M(u)^k\bigr),\quad \rho>0,
 \label{eq:logarithmic-normalization}
\end{equation}
where $B$ is a finite sum of positive-power logarithmic blocks, so that
$B(u)\to0$. The finite expression $A$ includes all retained nonvanishing
blocks, including any constant. Put $C(u)=\sigma e^{A(u)}$.

\begin{proposition}[Relative normalization]
Under these hypotheses,
\begin{equation}
 f(u)=C(u)\left(e^{B(u)}+
           \OO\bigl(u^\rho \M(u)^k\bigr)\right).
 \label{eq:relative-normalization}
\end{equation}
If $e^B=T+\OO(u^\eta \M(u)^\ell)$, the finite expression $C T$ has
absolute error
\begin{equation}
 \OO\left(|C(u)|
    \left(u^\rho \M(u)^k+u^\eta \M(u)^\ell\right)\right).
 \label{eq:normalized-absolute-error}
\end{equation}
\end{proposition}
\begin{proof}
The identity $f=C e^B e^R$ and the mean-value estimate
$|e^R-1|\le |R|e^{|R|}$ apply for real $R$.
Both $B$ and $R$ tend to zero, so $e^B$ and $e^{|R|}$ are bounded.
The first assertion follows, and the triangle inequality gives the second.
\end{proof}

The logarithmic remainder must be known to tend to zero; the hypothesis
$\rho>0$ ensures this. An error bound of order one in the logarithm
does not imply a relative error tending to zero in the original function. Likewise, a relative expansion of the logarithm
alone need not determine a relative expansion of its exponential.

\subsection{Products, varying powers, and coefficient extraction}

For positive real functions $f_j$ and real functions $r_j$, multiplication
and real exponentiation become addition before any truncation:
\begin{equation}
 \log\left(\prod_{j=1}^m f_j(u)^{r_j(u)}\right)
   =\sum_{j=1}^m r_j(u)\log f_j(u).
 \label{eq:product-logarithms}
\end{equation}
Cancellation can remove leading growth terms. It is therefore useful to
form the complete sum at the required precision before extracting $A$.
If $r_j(u)=\OO(u^{-s_j}\M(u)^{d_j})$, a logarithmic remainder
$\OO(u^h \M(u)^k)$ becomes
$\OO(u^{h-s_j}\M(u)^{k+d_j})$. To reach a desired vanishing absolute
error in the logarithm, the expansion of $\log f_j$ must be correspondingly
deeper. A real varying exponent cannot in general be treated as a constant
when propagating precision.

For an ordinary correction series
$B(t)=\sum_{n\ge1}a_n t^n$, write
$e^{B(t)}=\sum_{n\ge0}c_n t^n$. Formal differentiation gives
\begin{equation}
 c_0=1,\qquad
 c_n=\frac1n\sum_{j=1}^n j a_j c_{n-j}\quad(n\ge1).
 \label{eq:exponential-coefficient-recurrence}
\end{equation}
For real positive weights, the finite convolution algebra of
Section~\ref{sec:scale} supplies the same construction. If the smallest
positive weight is $\delta$, powers $B^n$ with $n\delta\ge h$ cannot
contribute below the strict weight cutoff $h$. Polynomial logarithmic
coefficients stay attached to their complete weight blocks.

\subsection{Stirling's logarithmic expansion}

On the positive real axis, Stirling's expansion is
\begin{equation}
 \log\Gamma(x)=
 \left(x-\tfrac12\right)\log x-x+\tfrac12\log(2\pi)
 +\sum_{n=1}^{N}
   \frac{B_{2n}}{2n(2n-1)x^{2n-1}}
 +\OO(x^{-2N-1}),\qquad x\to+\infty,
 \label{eq:stirling-logarithm}
\end{equation}
where $B_{2n}$ are Bernoulli numbers; see \cite[\S5.11]{dlmf}.
This is a Poincar\'e expansion: the assertion holds for every fixed
truncation order. It is not a convergence assertion as $N\to\infty$.
Multiplying, adding, and substituting into finitely many such expansions
is justified at each fixed order with the corresponding error estimates.

For products $\prod_j\Gamma(a_j(x))^{r_j(x)}$, assume that every argument
is eventually positive, tends to $+\infty$, and admits the substitutions
and error bounds needed in \eqref{eq:stirling-logarithm}. Formula
\eqref{eq:product-logarithms} then determines the logarithmic model.
Factors with arguments tending to positive finite points instead use
the local analytic expansion of $\log\Gamma$. Exponential normalization
combines the resulting growth and correction terms in either case.

\begin{example}[A square and a ratio]
Squaring the Gamma function doubles its logarithm. Exponentiating the
vanishing part gives
\begin{equation}
 \Gamma(x)^2=
 2\pi e^{-2x}x^{2x-1}
 \left(1+\frac1{6x}+\frac1{72x^2}
       -\frac{31}{6480x^3}-\frac{139}{155520x^4}
       +\OO(x^{-5})\right).
 \label{eq:gamma-square-expansion}
\end{equation}
For a ratio, subtract the logarithms first:
\begin{equation}
 \log\frac{\Gamma(3x)}{\Gamma(x)}
 = (3x-\tfrac12)\log3+2x\log x-2x
       -\frac1{18x}+\frac{13}{4860x^3}+\OO(x^{-5}).
\end{equation}
\begin{equation}
\begin{aligned}
 \frac{\Gamma(3x)}{\Gamma(x)}
 &=3^{3x-1/2}x^{2x}e^{-2x}
 \Bigl(1-\frac1{18x}+\frac1{648x^2}\\
 &\hspace{7em}+\frac{463}{174960x^3}-\frac{1867}{12597120x^4}
       +\OO(x^{-5})\Bigr).
\end{aligned}
 \label{eq:gamma-ratio-expansion}
\end{equation}
The remainder inside each parenthesis is relative to the prefactor.
The corresponding absolute error is that prefactor times $\OO(x^{-5})$.
\end{example}

\begin{example}[A varying exponent]
Multiplying \eqref{eq:stirling-logarithm} by $x$ turns the first
Stirling correction into the constant $1/12$. It therefore belongs to
the exact prefactor:
\begin{equation}
 \Gamma(x)^x=e^{1/12}
 \left(\sqrt{2\pi}\,e^{-x}x^{x-1/2}\right)^x
 \left(1-\frac1{360x^2}+\frac{1447}{1814400x^4}
       +\OO(x^{-6})\right).
 \label{eq:gamma-varying-power}
\end{equation}
The step from $\OO(x^{-7})$ in $\log\Gamma(x)$ to $\OO(x^{-6})$
in $x\log\Gamma(x)$ illustrates the required increase in working order.
\end{example}

\subsection{Barnes' \texorpdfstring{$G$}{G}-function and the argument shift}

Barnes' $G$-function satisfies
\begin{equation}
 G(z+1)=\Gamma(z)G(z),\qquad G(1)=1.
 \label{eq:barnes-recurrence}
\end{equation}
It is positive on the positive real axis. Write
$c_G=\zeta'(-1)=1/12-\log A$, where $A$ is Glaisher's constant.
For every fixed integer $N\ge0$, its logarithmic expansion is
\begin{equation}
\begin{aligned}
 \log G(z+1)
  ={}&\left(\frac{z^2}{2}-\frac1{12}\right)\log z
      -\frac{3z^2}{4}+\frac z2\log(2\pi)+c_G\\
    &+\sum_{k=1}^{N}
       \frac{B_{2k+2}}{2k(2k+2)z^{2k}}
       +\OO(z^{-2N-2}),\qquad z\to+\infty.
\end{aligned}
 \label{eq:barnes-logarithm}
\end{equation}
This form follows by inserting Stirling's expansion for
$\log\Gamma(z+1)$ into \cite[(5.17.5)]{dlmf}; the relation between
$A$ and $\zeta'(-1)$ is \cite[(5.17.7)]{dlmf}. The factor $z$ multiplying
that logarithm requires one additional order of its absolute error.
Expanding it through the Bernoulli term with index $2N+2$ gives the
remainder displayed in \eqref{eq:barnes-logarithm}. As with Stirling's
formula, each fixed truncation has its stated error, without asserting
convergence of the infinite series.

The first vanishing logarithmic terms are
\begin{equation}
 -\frac1{240z^2}+\frac1{1008z^4}
 -\frac1{1440z^6}+\frac1{1056z^8}+\OO(z^{-10}).
 \label{eq:barnes-logarithmic-corrections}
\end{equation}
Define the positive prefactor
\begin{equation}
 P_G^+(z)=e^{c_G}(2\pi)^{z/2}
          e^{-3z^2/4}z^{z^2/2-1/12}.
 \label{eq:barnes-shifted-prefactor}
\end{equation}
Relative normalization and
\eqref{eq:exponential-coefficient-recurrence} give, for example,
\begin{equation}
 G(z+1)=P_G^+(z)
 \left(1-\frac1{240z^2}+\frac{269}{268800z^4}
       +\OO(z^{-6})\right).
 \label{eq:barnes-shifted-expansion}
\end{equation}
Every correction power in this coordinate is even. A logarithmic
remainder $\OO(z^{-2N-2})$ becomes a relative remainder of that
same order after exponentiation; the absolute remainder includes
the positive factor $P_G^+(z)$.

For $G(x)$ the argument in \eqref{eq:barnes-logarithm} is
$z=x-1$. Retaining powers of $1/(x-1)$ gives the shifted form above.
Expressing the result in powers of $1/x$ changes both the exact
nonvanishing logarithmic part and the correction coefficients. The
recurrence \eqref{eq:barnes-recurrence} provides an equivalent direct
derivation: subtract \eqref{eq:stirling-logarithm} from the expansion
of $\log G(x+1)$. Put
\begin{equation}
 P_G(x)=e^{c_G}(2\pi)^{(x-1)/2}
        e^{-3x^2/4+x}x^{x^2/2-x+5/12}.
 \label{eq:barnes-prefactor}
\end{equation}
Then
\begin{equation}
\begin{aligned}
 \log G(x)={}&\log P_G(x)-\frac1{12x}-\frac1{240x^2}\\
            &+\frac1{360x^3}+\frac1{1008x^4}+\OO(x^{-5}),
\end{aligned}
 \label{eq:barnes-unshifted-logarithm}
\end{equation}
and its first five relative blocks are
\begin{equation}
\begin{aligned}
 G(x)=P_G(x)\Bigl(&1-\frac1{12x}-\frac1{1440x^2}
              +\frac{157}{51840x^3}\\
              &+\frac{65911}{87091200x^4}+\OO(x^{-5})\Bigr).
\end{aligned}
 \label{eq:barnes-five-blocks}
\end{equation}
The term $z=x-1$ is therefore essential when translating the
logarithmic formula to an expansion of $G(x)$. A shift of its argument
does not preserve the correction bracket in a fixed coordinate.

Products of positive Gamma and Barnes factors use
\eqref{eq:product-logarithms}, so their logarithmic expansions can
be combined before choosing the prefactor. Real varying powers
transport the logarithmic error with the precision loss described
above. Positive finite limiting arguments instead use local analytic
logarithms. Integer shifts can also produce exact simplifications;
for example,
\[
 \frac{G(x+1)}{G(x)}=\Gamma(x),\qquad
 \frac{G(x+1)}{\Gamma(x)G(x)}=1\qquad(x>0).
\]
These identities justify exact cancellation independently of a finite
asymptotic coefficient calculation.

\subsection{Exact identities and elementary growth}

The identities $\Gamma(x+1)=x\Gamma(x)$ and
$\mathrm B(a,b)=\Gamma(a)\Gamma(b)/\Gamma(a+b)$ connect the same
analysis to factorials, binomial coefficients, beta functions, and
rising factorials. For example,
\begin{equation}
 \binom{2x}{x}=\frac{4^x}{\sqrt{\pi x}}
 \left(1-\frac1{8x}+\frac1{128x^2}
       +\frac5{1024x^3}+\OO(x^{-4})\right)
 \qquad(x\to+\infty),
\end{equation}
where the binomial coefficient is defined by its Gamma quotient.
The recurrence also proves $\Gamma(x+1)/\Gamma(x)=x$ exactly.
Vanishing coefficients in a finite asymptotic computation alone would
not prove this identity: a function may have all algebraic coefficients
zero and still have a nonzero flat remainder.

The construction applies equally to elementary growth. Two examples are
\begin{align}
 e^{x+1/x}
   &=e^x\left(1+\frac1x+\frac1{2x^2}+\OO(x^{-3})\right),\\
 x^{x+1/x}
   &=x^x\left(1+\frac{\log x}{x}
             +\frac{(\log x)^2}{2x^2}
             +\OO\left(\frac{(\log x)^3}{x^3}\right)\right).
\end{align}
In each case it is the vanishing logarithmic correction, rather than
the entire growing function, that is expanded in a small scale.
